% s14_t3_maj: the two-loop majorants at every photon mass: torus floor, monotonicity, closed form (input from sec/s14.tex)

\subsection*{The two-loop orders at every photon mass}

From here to the end of the proof of \cref{prop:s14:T3} we use the units $m=1$ ($\lambda$ and $e^2$ in units of $m$, $\sigma$ in units of $m^2$), $N=2$ and $M_1=0.7$; ``$\eps_k\le1$'' means $\eps_k\le1/m$. The two-loop coefficients $S_{21,k}$, $S_{40,k}$ are bounded from below through the torus paper's Appendix~G \cite{SymPhase}, whose statements are applied with their hypotheses matched here; that appendix is written on periodic tori $T$ of side $\ell_T\ge\ell_0$ with $\ell_T/\eps\in2\N$, in units $m=1$, and its statements carry the reference photon mass $\mu_*\in(0,1]$ of its photon propagator $\hat D=(\hat k^2+\mu_*^2)^{-1}$ as a parameter; we read them at $\mu_*=\lmu$. Its objects are the following.

\begin{enumerate}[label=\textup{(G\arabic*)},leftmargin=*]
\item \emph{The torus coefficient} $a_{K,T}(M)$ \cite[Def.~2.2]{SymPhase}: King's recursion \eqref{eq:s14:CTdef} computed on $T$ with the periodic propagators of $T$ at scalar mass $M$ and photon mass $\mu_*$. Read at $\mu_*=\lmu$, this is the bare coefficient $\alpha^{\mathrm{tor}}_{k,T}(M)=M^2+\dm_{k,T}(M)$ of King's torus convention (\cref{def:s10:torus}(ii)): both are the same recursion with the same propagators. Its exact two-loop form \cite[Lemma~G.1]{SymPhase} is, at $N=2$,
\[
a_{K,T}(M)=M^2-16\lambda C_T(M)-2e^2D_T-e^2J_T(M)+128\lambda^2S_T(M)+16e^2\lambda F_T(M)+e^4Q_T(M),
\]
with $C_T(M)=\int_T\hat G_M$, $D_T=\int_T\hat D$, $J_T(M)=M^2\int_T\hat D\hat G_M$, $S_T(M)=\sum_{x\in T}\eps^3G_M(x)^3$, $F_T$ the order-$e^2\lambda$ and $Q_T$ the order-$e^4$ coefficient, the term $\frac34\eps^2D_T^2$ of $Q_T$ being independent of $M$; $M\mapsto a_{K,T}(M)$ is smooth. (The three low orders are \cref{lem:s14:low} on the torus.)
\item \emph{The one-loop majorants} \cite[Lemma~G.15]{SymPhase}: for $\alpha,\beta\ge0$ with $g\defeq\alpha+\beta>\frac32$, $a,b>0$, $\nu\defeq\min(a,b)$ and $\ell\in(0,\infty]$,
\begin{equation}\label{eq:s14:Iab}
\mathcal I_{\alpha,\beta}(a,b;\ell)\defeq c_NC_3+\frac{3\pi}{16\ell}C_2+\frac3{4\ell^2}C_1+\frac{a^{-2\alpha}b^{-2\beta}}{\ell^3},\qquad c_N=\frac{\pi^3}8,
\end{equation}
each $C_d=C_d(\alpha,\beta;a,b)$ the minimum of its admissible closed forms: for $C_3$ the $\nu$-form $\Gamma(g-\frac32)\nu^{3-2g}/((4\pi)^{3/2}\Gamma(g))$, the Beta form in $a$, $\frac1{4\pi^2}B(c,\alpha-c)a^{2c-2\alpha}$ with $c=\frac{3-2\beta}2$ if $\beta<\frac32$, the same with $(a,\alpha)\leftrightarrow(b,\beta)$ if $\alpha<\frac32$, the exact values $\frac1{4\pi(a+b)}$ at $(\alpha,\beta)=(1,1)$ and $\frac1{8\pi a(a+b)^2}$ at $(2,1)$, and the H\"older form $\gamma_3(g)a^{(3-2g)\alpha/g}b^{(3-2g)\beta/g}$, $\gamma_3(g)a^{3-2g}$ the closed form of $\int(u+a^2)^{-g}$; for $C_2$ the $\nu$-form $\nu^{2-2g}/(g-1)$, the Beta form $a^{2(1-g)}B(1-\beta,g-1)$ if $0<\beta<1$ and its symmetric form, and the H\"older form; for $C_1$ the $\nu$-form $\sqrt\pi\Gamma(g-\frac12)\nu^{1-2g}/\Gamma(g)$, the Beta form $a^{1-2g}B(\frac12-\beta,g-\frac12)$ if $0<\beta<\frac12$ and its symmetric form, and the H\"older form. For every torus of side $\ell$ and every $\eps$, $\int_T(\hat k^2+a^2)^{-\alpha}(\hat k^2+b^2)^{-\beta}\le\mathcal I_{\alpha,\beta}(a,b;\ell)$, on $\eps\Z^3$ with $\ell=\infty$; $\mathcal I_{\alpha,\beta}$ is decreasing in $\ell$, and at $(a,b)=(M,\mu)$ every admissible form is non-increasing in $M$. The majorants, all at $(M,\mu;\ell_0)$ with $\mu=\mu_*$, are $\bar\tau=\mathcal I_{2,0}$, $\bar\psi=\mathcal I_{1,1}$, $\bar\delta=\mathcal I_{0,2}$, $\bar\Theta=\mathcal I_{1,3/4}+\mathcal I_{3/4,1}$, $\bar J_s=\mathcal I_{3-s/2,0}$ and $\bar J'_s=\mathcal I_{2-s/2,1}$ for $s\in\{0,1,\frac32\}$, $\bar\Lambda_0=\mathcal I_{3/4,1}$, $\bar\Lambda_1=\mathcal I_{5/4,1}+\mathcal I_{2,1/4}$, $\bar\Lambda_2=\mathcal I_{1/4,2}+\mathcal I_{1,5/4}$: fourteen distinct pairs $(\alpha,\beta)$, all with $g\ge\frac74>\frac32$; and $\bar\omega\defeq w+\hat E(\mu,\ell_0)$ with $w=0.252731$ (used as $w\le0.25273101$), $\hat E\defeq\min\{\bar E,\tilde E\}$, where $\bar E(\mu,\ell)=\frac{7.37}\ell+\frac{3\pi}{4\ell}E_1(\frac{\mu^2\ell^2}8)+\frac{12\pi\,\mathrm{erfc}(\mu\ell/(2\sqrt2))}{\mu\ell^2}+\frac{64e^{-\mu^2\ell^2/8}}{\mu^2\ell^3}$ \cite[Lemma~H.31]{SymPhase} and $\tilde E\ge0$ \cite[Lemma~G.14]{SymPhase}; for $\eps\le1$ and $\ell_T\ge\ell_0$, $\eps D_T\le\bar\omega$ \cite[Lemma~G.14(c)]{SymPhase}, and on $\eps\Z^3$ ($\ell_0=\infty$) one takes $\bar\omega=w$ \cite[Lemma~G.24]{SymPhase}.
\item \emph{The slope bounds} \cite[Lemma~G.2]{SymPhase}: for $\eps\le1$, $\ell_T\ge\ell_0$, $M>0$, with the majorants at $(M,\mu_*;\ell_0)$, $\bar\beta\defeq2M[2\bar\tau+8\bar\psi]$ and
\begin{gather*}
j\defeq2M\min\{\bar\tau,\bar\psi\},\qquad \varphi_3\defeq4M\big[2\sqrt2\,\bar\Theta\bar J_{3/2}+|\mu_*^2-2M^2|\bar\psi\bar J_0+3\sqrt3\,\bar\omega\bar J_1\big]+\bar\beta\bar\tau,\\
b_4\defeq\varphi_A+28M\bar\delta\bar\tau+\varphi_{D2},\quad \varphi_A\defeq6M\big[2\sqrt2\,\bar\Theta\bar J'_{3/2}+|\mu_*^2-2M^2|\bar\psi\bar J'_0+3\sqrt3\,\bar\omega\bar J'_1\big]+\bar\beta\bar\psi,\\
\varphi_{D2}\defeq2M\bar\psi(2\bar\psi+4\bar\delta)+\min\{8M\bar\psi(6\bar\psi+2\bar\tau+4\bar\delta),\ 2^{5/2}M(2\bar\Lambda_1+4\bar\Lambda_2)\bar\Lambda_0\},
\end{gather*}
one has $0\le\partial_MJ_T\le j(M)$, $0\le-128\,\partial_MS_T\le768M\bar\tau(M)^2$, $|\partial_MF_T|\le\varphi_3(M)$ and $\partial_MQ_T\ge-b_4(M)$.
\item \emph{The two-loop gap} \cite[Prop.~G.3]{SymPhase}: for $0<M_1<M_2$, $\mu_*\in(0,1]$, $\ell_0>0$, every $\eps\le1$ and every torus of side $\ell_T\ge\ell_0$, $a_{K,T}(M_2)-a_{K,T}(M_1)\ge(M_2^2-M_1^2)+16\lambda h_1-e^2\bar j-\lambda^2\bar s-e^2\lambda\bar B_3-e^4\bar B_4^-$ with $h_1\ge0$, $\bar j=\int_{M_1}^{M_2}j$, $\bar s=\int768M\bar\tau^2$, $\bar B_3=16\int\varphi_3$, $\bar B_4^-=\int b_4$; its proof writes each order as $\int_{M_1}^{M_2}\partial_M(\cdot)$ and applies (G3) \emph{order by order}: $-[J_T(M_2)-J_T(M_1)]\ge-\bar j$, $128[S_T(M_2)-S_T(M_1)]\ge-\bar s$, $16[F_T(M_2)-F_T(M_1)]\ge-\bar B_3$, $Q_T(M_2)-Q_T(M_1)\ge-\bar B_4^-$.
\item \emph{The cell rule} (the proof of \cite[Cor.~G.4]{SymPhase}): for a density $\phi\in\{j,768M\bar\tau^2,16\varphi_3,b_4\}$ and a cell $[a,b]\subset(0,\infty)$ let $\phi^{\rm cell}[a,b](\mu,\ell_0)$ be the number obtained from the formula for $\phi$ by evaluating every majorant at the left endpoint $a$, the prefactor $M$ at the right endpoint $b$, the prefactor $|\mu^2-2M^2|$ at its maximum over $\{a,b\}$ and $\bar\omega$ as it is; since the majorants are non-negative and non-increasing in $M$ and the formulas are sums, products and minima with non-negative coefficients, $\phi\le\phi^{\rm cell}[a,b]$ on $[a,b]$, and for a partition $M_1=M^{(0)}<\dots<M^{(n)}=M_2$
\begin{equation}\label{eq:s14:cell}
\int_{M_1}^{M_2}\phi\,dM\ \le\ \Phi_n(\mu,\ell_0)\defeq\sum_{i=0}^{n-1}\big(M^{(i+1)}-M^{(i)}\big)\,\phi^{\rm cell}[M^{(i)},M^{(i+1)}](\mu,\ell_0).
\end{equation}
We write $\bar j_n,\bar s_n,\bar B_{3,n},\bar B^-_{4,n}$ for the four upper sums on the uniform partition of $[0.7,1]$ into $n$ cells, and $\bar X_n$ for any of them; $n=400$ unless said otherwise.
\end{enumerate}

\begin{lemma}[the four orders on $\Lk$ at every photon mass]\label{lem:s14:E}
Let $\lmu\in(0,1]$, $\ell_0>0$, $n\ge1$. For every $k$ with $\eps_k\le1$:
\[
2S_{20,k}\ge-\bar j_n(\lmu,\ell_0),\qquad 2S_{02,k}\ge-\bar s_n,\qquad 2S_{21,k}\ge-\bar B_{3,n}(\lmu,\ell_0),\qquad 2S_{40,k}\ge-\bar B^-_{4,n}(\lmu,\ell_0).
\]
\end{lemma}

\begin{proof}
Fix such a $k$, and let $T=T_J$ be King's tori (\cref{def:s10:torus}; side $2MJ$, $2MJ\cdot5^k$ sites per side, an even number), $J$ so large that $2MJ\ge\ell_0$. By (G1) at $\mu_*=\lmu$, the coefficients of the torus-convention bare coefficient $\alpha^{\mathrm{tor}}_{k,T}(M)$ in the five orders are the terms of the exact two-loop form, so that for $S^T_{\alpha\beta,k}\defeq\frac12[\dm_{(\alpha\beta),T}(\eps_k;1)-\dm_{(\alpha\beta),T}(\eps_k;0.7)]$,
\[
2S^T_{20,k}=-[J_T(1)-J_T(0.7)],\quad 2S^T_{02,k}=128[S_T(1)-S_T(0.7)],\quad 2S^T_{21,k}=16[F_T(1)-F_T(0.7)],\quad 2S^T_{40,k}=Q_T(1)-Q_T(0.7),
\]
the $M$-independent terms $-2D_T$ and $\frac34\eps^2D_T^2$ cancelling in the differences. The hypotheses of (G3) and (G4) are met: $\eps_k\le1$, $\ell_T\ge\ell_0$, $\mu_*=\lmu\in(0,1]$, $M\in[0.7,1]$, and the lemmas behind (G3) carry $\mu_*$ as a parameter in their statements ((G2): $a,b>0$ and $g>\frac32$ for every instance; $\eps\le1$ and $\ell_T\ge\ell_0$ for $\bar\omega$; $\ell_T/\eps\in\N$, $\ell_T/\eps\ge2$ for $\bar E$). Hence, by the order-by-order bounds of (G4) with the majorants at $(\lmu,\ell_0)$ and by the cell rule \eqref{eq:s14:cell}, $2S^T_{20,k}\ge-\bar j_n(\lmu,\ell_0)$, $2S^T_{02,k}\ge-\bar s_n$, $2S^T_{21,k}\ge-\bar B_{3,n}(\lmu,\ell_0)$, $2S^T_{40,k}\ge-\bar B^-_{4,n}(\lmu,\ell_0)$, uniformly in $J$ ($\bar s$ does not depend on $\mu$: $\bar\tau=\mathcal I_{2,0}$ has $\beta=0$). Now let $J\to\infty$ at fixed $k$. By \cref{lem:s10:V}, $\alpha^{\mathrm{tor}}_{k,T_J}(M)\to\alpha^\infty_k(M)$ for $M\in\{0.7,1\}$, and its proof (\cref{ssec:a06:V}) establishes the convergence of the recursion \eqref{eq:s14:CTdef} order by order: $\dm_{(\alpha\beta),T_J}(\eps_k;M)\to\dm_{(\alpha\beta)}(\eps_k;M)$ for each of the five orders (alternatively: a polynomial of bounded degree in $(e^2,\lambda)$ that converges at every point of an open set converges coefficient by coefficient, and \cref{lem:s10:V} holds at every point of King's range). So $S^{T_J}_{\alpha\beta,k}\to S_{\alpha\beta,k}$, and closed inequalities pass to the limit.
\end{proof}

\begin{lemma}[the torus floor can be removed]\label{lem:s14:TF}
For every $\lmu\in(0,1]$, $n\ge1$ and every $k$ with $\eps_k\le1$, the four inequalities of \cref{lem:s14:E} hold with $\bar X_n(\lmu,\infty)$, the upper sums formed with the $\eps\Z^3$ majorants $\mathcal I_{\alpha,\beta}(M,\lmu;\infty)=c_NC_3(M,\lmu)$ and with $\bar\omega=w$.
\end{lemma}

\begin{proof}
Fix $k$, $\lmu$, $n$. By \cref{lem:s14:E}, $2S_{\alpha\beta,k}\ge-\bar X_n(\lmu,\ell_0)$ for \emph{every} $\ell_0>0$: the torus floor is a free parameter of (G4), and the tori $T_J$ of the volume limit eventually exceed any $\ell_0$. It therefore suffices that $\bar X_n(\lmu,\ell_0)\to\bar X_n(\lmu,\infty)$ as $\ell_0\to\infty$. Each $\bar X_n$ is a finite sum of cell values, and each cell value is a fixed finite expression (sums, products and minima, with the fixed prefactors) in the majorants $\mathcal I_{\alpha,\beta}(a,\lmu;\ell_0)$ at the left endpoint $a$ of the cell and in $\bar\omega(\lmu,\ell_0)$, hence continuous in these. As $\ell_0\to\infty$, $\mathcal I_{\alpha,\beta}(a,\lmu;\ell_0)\to c_NC_3(a,\lmu)$ by \eqref{eq:s14:Iab}, since $C_2$ and $C_1$ are finite numbers that do not depend on $\ell_0$; and $0\le\hat E(\lmu,\ell_0)\le\bar E(\lmu,\ell_0)\to0$, each of the four terms of $\bar E$ being a positive constant times $\ell_0^{-1}$, $\ell_0^{-1}E_1(\lmu^2\ell_0^2/8)$, $\ell_0^{-2}\mathrm{erfc}(\lmu\ell_0/(2\sqrt2))$ or $\ell_0^{-3}e^{-\lmu^2\ell_0^2/8}$, so $\bar\omega\to w$.
\end{proof}

\begin{lemma}[monotonicity in the photon mass; interval bounds]\label{lem:s14:MON}
Let $0<\mu_{\rm lo}\le\mu_{\rm hi}\le1$. Define the \emph{interval upper sum} $\bar X^{\rm int}_n([\mu_{\rm lo},\mu_{\rm hi}])$ as $\bar X_n(\cdot,\infty)$ with every majorant evaluated at $\mu=\mu_{\rm lo}$ and, on the cell $[a,b]$, the prefactor $|\mu^2-2M^2|$ replaced by $\max\{|\mu^2-2M^2|:\mu\in\{\mu_{\rm lo},\mu_{\rm hi}\},\,M\in\{a,b\}\}$. Then $\bar X_n(\mu,\infty)\le\bar X^{\rm int}_n([\mu_{\rm lo},\mu_{\rm hi}])$ for every $\mu\in[\mu_{\rm lo},\mu_{\rm hi}]$, so the inequalities of \cref{lem:s14:TF} hold with $\bar X^{\rm int}_n([\mu_{\rm lo},\mu_{\rm hi}])$ for all $\lmu\in[\mu_{\rm lo},\mu_{\rm hi}]$ at once. Moreover, if $\mu_{\rm hi}\le0.989949$, the interval upper sum equals the point upper sum $\bar X_n(\mu_{\rm lo},\infty)$.
\end{lemma}

\begin{proof}
A cell value is non-decreasing in each majorant and in the prefactor (sums, products and minima of non-negative quantities with non-negative coefficients). (i) Each admissible closed form of $C_3$ is non-increasing in $b=\mu$: the $\nu$-form, since $\nu=\min(a,b)$ is non-decreasing in $b$ and $3-2g<0$; the Beta form in $a$ is $b$-free; the Beta form in $b$ carries $b^{2c'-2\beta}$ with $c'=\frac{3-2\alpha}2<\beta$, as $g>\frac32$; the H\"older form carries $b^{(3-2g)\beta/g}$; the exact forms $\frac1{4\pi(a+b)}$ and $\frac1{8\pi a(a+b)^2}$ decrease in $b$. A minimum of non-increasing functions is non-increasing, so every $\mathcal I_{\alpha,\beta}(a,\mu;\infty)$ is non-increasing in $\mu$, and $\bar\omega=w$ does not depend on $\mu$. (ii) $\mu\mapsto|\mu^2-2M^2|$ is convex in $\mu^2$ and $M\mapsto|\mu^2-2M^2|$ convex in $M^2$, so the maximum over the rectangle $[\mu_{\rm lo},\mu_{\rm hi}]\times[a,b]$ is attained at a corner. Hence every ingredient of the cell value at $\mu$ is at most the corresponding ingredient of the interval cell value. For the last claim: on $[0.7,1]$ one has $2M^2\ge0.98>0.989949^2$, so $|\mu^2-2M^2|=2M^2-\mu^2$ on the whole rectangle, maximal at $(\mu_{\rm lo},b)$, which is also the point prefactor at $\mu_{\rm lo}$. (For the interval $[0.98,0.99]$ the corner $(0.99,0.7)$ of the first cell gives $|0.9801-0.98|=10^{-4}$, below the value $2\cdot0.70075^2-0.98^2=0.0217$ at $(0.98,0.70075)$, so the equality holds there too; only the interval $[0.99,1]$ differs, by $8.5786210$ against the point value $8.5785880$ of $\bar B_3$.)
\end{proof}

\begin{lemma}[a closed form, uniform in the photon mass]\label{lem:s14:CF}
For $M\in[0.7,1]$ put $\bar\tau^*(M)\defeq\mathcal I_{2,0}(M;\infty)=\frac{c_N}{8\pi M}$, $\bar\psi^*(M)\defeq\frac{c_N}{4\pi M}$, and for $X\in\{\Theta,J_s,J'_s\}$ let $X^*(M)$ be $c_N$ times the minimum over the $\mu$-free admissible forms of $C_3$ only (the Beta forms in $a=M$ and, for $\beta=0$, the H\"older form $\gamma_3(g)M^{3-2g}$; the $\nu$-forms, the Beta forms in $b$, the exact form at $(2,1)$ and the H\"older forms with $\beta>0$ are dropped, and the exact form at $(1,1)$ is replaced by its bound $\frac1{4\pi a}$, which is $\bar\psi^*$). Put $\bar\beta^*\defeq2M[2\bar\tau^*+8\bar\psi^*]$,
\begin{gather*}
\varphi_3^*(M)\defeq4M\big[2\sqrt2\,\bar\Theta^*\bar J^*_{3/2}+2M^2\bar\psi^*\bar J^*_0+3\sqrt3\,w\,\bar J^*_1\big]+\bar\beta^*\bar\tau^*,\qquad
\varphi_A^*(M)\defeq6M\big[2\sqrt2\,\bar\Theta^*\bar J'^*_{3/2}+2M^2\bar\psi^*\bar J'^*_0+3\sqrt3\,w\,\bar J'^*_1\big]+\bar\beta^*\bar\psi^*,\\
P(M)\defeq\varphi_A^*+52M\bar\psi^{*2}+16M\bar\psi^*\bar\tau^*,\qquad Q(M)\defeq\frac{c_N}{8\pi}M\big(28\bar\tau^*+40\bar\psi^*\big).
\end{gather*}
Then $16\varphi_3(M;\mu,\infty)\le16\varphi_3^*(M)$ and $b_4(M;\mu,\infty)\le P(M)+Q(M)/\mu$ for all $M\in[0.7,1]$ and $\mu\in(0,1]$. Consequently, with the $400$-cell upper sums $B_3^*\defeq\sum h\,(16\varphi_3^*)^{\rm cell}$, $P^*\defeq\sum hP^{\rm cell}$, $Q^*\defeq\sum hQ^{\rm cell}$ (majorants at the left endpoint, the prefactors $M$ and $2M^2$ at the right endpoint), for every $\lmu\in(0,1]$ and every $k$ with $\eps_k\le1$:
\[
2S_{21,k}\ge-B_3^*,\qquad 2S_{40,k}\ge-\big(P^*+Q^*/\lmu\big),\qquad B_3^*=13.187549,\quad P^*=4.7629000,\quad Q^*=0.77120851 .
\]
\end{lemma}

\begin{proof}
Each starred majorant is an upper bound of the corresponding majorant at $(M,\mu;\infty)$ for every $\mu>0$: a minimum over a subset of the admissible forms is at least the minimum over all of them, and $\mathcal I_{1,1}(M,\mu;\infty)\le c_N\frac1{4\pi(M+\mu)}\le\frac{c_N}{4\pi M}$. Further $\bar\delta(\mu)=\mathcal I_{0,2}(M,\mu;\infty)\le c_N\frac1{4\pi^2}B(\frac32,\frac12)\mu^{-1}=\frac{c_N}{8\pi\mu}$ (the Beta form in $b$, admissible since $\alpha=0<\frac32$ and $\beta=2>\frac32$); $|\mu^2-2M^2|\le2M^2$, since $\mu^2\le1\le4M^2$; $\bar\omega=w$ (\cref{lem:s14:TF}); and in $\varphi_{D2}$ the minimum is at most its first branch. Substituting these bounds into the formulas (G3), which are non-decreasing in every majorant, gives $16\varphi_3\le16\varphi_3^*$ and
\[
b_4\le\varphi_A^*+28M\bar\tau^*\frac{c_N}{8\pi\mu}+2M\bar\psi^*\Big(2\bar\psi^*+\frac{4c_N}{8\pi\mu}\Big)+8M\bar\psi^*\Big(6\bar\psi^*+2\bar\tau^*+\frac{4c_N}{8\pi\mu}\Big)=P(M)+\frac{Q(M)}\mu .
\]
The starred majorants are non-increasing in $M$ (powers $M^{2c-2\alpha}$ with $c<\alpha$, and $1/M$), so the cell rule \eqref{eq:s14:cell} applies to $16\varphi_3^*$, $P$ and $Q$ with the prefactors $M$, $2M^2$ at the right endpoint. The two inequalities follow \emph{cellwise} from \cref{lem:s14:TF}: for a cell $[a,b]$ the cell value of $16\varphi_3$ (resp.\ $b_4$) at $(\lmu,\infty)$ is a non-decreasing expression in the majorants at $(a,\lmu;\infty)$, in $\bar\omega$ and in the prefactor, each of which is at most its starred counterpart at $a$ (with $\bar\delta\le c_N/(8\pi\lmu)$, $\max_{M\in\{a,b\}}|\lmu^2-2M^2|\le2b^2$, $\bar\omega=w$); so the cell value is at most the cell value of $16\varphi_3^*$ (resp.\ of $P+Q/\lmu$). Summing over the $400$ cells, $\bar B_{3,400}(\lmu,\infty)\le B_3^*$ and $\bar B^-_{4,400}(\lmu,\infty)\le P^*+Q^*/\lmu$, and \cref{lem:s14:TF} gives the two inequalities. The three numbers are the $400$-cell sums, computed as described after \cref{tab:s14:T3}. (The slope bounds (G3) are not invoked with the starred majorants in place of the one-loop majorants: as printed they make only $\bar\tau$, $\bar\psi$ generic; the domination above is all that is used.)
\end{proof}
